\begin{proposition}[G-13: An explicit finite bound for the first rectangular gap]\label{G-13}
For $n\geq4$ put
\[
 \beta=6\binom{n+3}{3}n^{30},\qquad u=2n^2-13,\qquad
 C=(n\beta^2)^u,
\]
and define
\begin{equation}\label{G13:first-bound}
 s_*=(u+1)(C-1)+1,\qquad d_* = \beta s_*.
\end{equation}
There is an integer $d$ with $1\leq d\leq d_*$ and $4\mid nd$
such that, for $\lambda=(nd/4)^4$,
\[
 m_\lambda(\operatorname{per}_n)>m_\lambda(\det_n).
\]
The same conclusion holds with $\operatorname{per}_n$ replaced by
$\ell^{n-m}\operatorname{per}_m$ whenever $n>m\geq4$ and
\[
 \dim\mathcal P_{m,4}+3>2n^2+2,
 \qquad n<4m/3.
\]
In particular it applies to $n=m+1$, $m\geq8$, and, by R-18,
to $n=\lfloor\rho m\rfloor$ for all sufficiently large $m$
and each fixed $1<\rho<\sqrt{3/2}$.

The bound satisfies $\log d_*=O(n^2\log n)$. This is an
existence bound within a finite interval. It does not specify
the individual successful degree or replace G-11's assertion
at every degree of its displayed progression.
\end{proposition}

\paragraph{Background and scope.}
We use the stable points, exact invariant quotients, and quotient
dimensions established in G-11, together with Derksen (2001),
Theorem~1.1 and Example~2.1, printed p.~956. We also use the
projective image-degree argument of G-10 and Chardin's
unnumbered main theorem (1989), printed p.~306. The argument
below compares a weighted degree filtration with the ordinary
degree filtration in an affine generator embedding; it needs
no factorial Veronese.

\begin{proof}
Let $G=\mathrm{SL}_4(\mathbb C)$ and let $V_n$ be the space
of quaternary forms of degree $n$. Write
\[
 A_D=\mathbb C[\mathcal D_{n,4}]^G,
 \qquad A_X=\mathbb C[X]^G,
\]
where $X$ is $\mathcal P_{n,4}$ or the padded cone in the
statement. The stable-point and dimension arguments of G-11
give
\begin{equation}\label{G13:dims}
 t:=\dim A_D\leq u,\qquad \dim A_X\geq u+1.
\end{equation}
These are positively graded domains, with degree-zero part
$\mathbb C$, and both are graded quotients of
$\mathbb C[V_n]^G$. Choose homogeneous generators
$g_1,\ldots,g_r$ of the latter ring with positive degrees
$e_i\leq\beta$. Their images generate both quotient rings.

First bound the cumulative graded dimension of $A_D$. The
images of the generators give a closed affine embedding
\[
 Y_D=\operatorname{Spec} A_D\lhook\joinrel\longrightarrow
 \mathbb A^r.
\]
Compose this embedding with the determinant-pencil map and
the affine quotient. The resulting map from
$\mathbb A^{4n^2}$ has coordinates
$g_i(\det L)$, homogeneous of degrees $ne_i\leq n\beta$
in the pencil parameters. Its image is dense in $Y_D$.
To homogenize it, adjoin a source coordinate $z$ and use
the degree-$n\beta$ coordinate polynomials
\[
 z^{n\beta},\qquad
 z^{n\beta-ne_i}g_i(\det L),\quad 1\leq i\leq r.
\]
They define a rational map from $\mathbb P^{4n^2}$ whose
image closure is the ordinary projective closure
$\overline{Y_D}\subseteq\mathbb P^r$. It has dimension
$t$. G-10's image-degree bound, allowing a base locus, gives
\begin{equation}\label{G13:affine-degree}
 \deg\overline{Y_D}\leq(n\beta)^t\leq(n\beta)^u.
\end{equation}

For $d\geq0$, let
\[
 H_D^{\leq}(d)=\sum_{j=0}^d\dim(A_D)_j.
\]
Every element of $(A_D)_j$ has a preimage homogeneous of
weighted degree $j$ in $\mathbb C[T_1,\ldots,T_r]$ with
$\deg T_i=e_i$. Since every $e_i\geq1$, each monomial in
that preimage has ordinary degree at most $j$. It follows
that the direct sum of the graded pieces through $d$ is
contained in the space of functions on $Y_D$ represented
by ordinary polynomials of degree at most $d$.
This is an inclusion of filtered vector spaces, not an
identification of the two gradings.

The latter ordinary filtered dimension equals the
degree-$d$ Hilbert function of $\overline{Y_D}$:
homogenization with its first coordinate, followed by
dehomogenization on the dense affine chart, identifies the
two quotients. The projective closure is integral over
$\mathbb C$. Also $t\geq19$ by the fixed block-padding and
stable-quotient observation in G-11, so the positive-dimension
hypothesis in Chardin's stated theorem holds. That theorem
and \eqref{G13:affine-degree} yield
\begin{equation}\label{G13:upper}
 H_D^{\leq}(d)
 \leq(n\beta)^t\binom{d+t}{t}
 \leq(n\beta)^u\binom{d+u}{u}.
\end{equation}

For the other ring, a maximal algebraically independent
subset of the images of $g_1,\ldots,g_r$ has size
$\dim A_X$. Indeed those images generate its fraction
field, and maximality makes every remaining generator
algebraic over the chosen subset. By \eqref{G13:dims},
choose $u+1$ algebraically independent images. Each has
degree at most $\beta$. Their monomials of ordinary total
degree at most $s$ are linearly independent and have
graded degree at most $\beta s$. Consequently
\begin{equation}\label{G13:lower}
 H_X^{\leq}(\beta s)
 :=\sum_{j=0}^{\beta s}\dim(A_X)_j
 \geq\binom{s+u+1}{u+1}.
\end{equation}

For every nonnegative integer $s$, the exact binomial ratio
satisfies
\begin{align*}
 \frac{\binom{s+u+1}{u+1}}{\binom{\beta s+u}{u}}
 &=\frac{s+u+1}{u+1}
     \prod_{i=1}^u\frac{s+i}{\beta s+i}\\
 &\geq\frac{s+u+1}{(u+1)\beta^u}.
\end{align*}
The last inequality follows term by term from
$\beta(s+i)\geq\beta s+i$. At $s=s_*$ one has
\[
 s_*+u+1=(u+1)C+1,
\]
so the final lower bound is strictly greater than
$C/\beta^u=(n\beta)^u$. Equations
\eqref{G13:upper}--\eqref{G13:lower} therefore give
\[
 H_X^{\leq}(d_*)>H_D^{\leq}(d_*).
\]
The degree-zero dimensions are both one. Hence for some
$1\leq d\leq d_*$ the individual graded pieces satisfy
$\dim(A_X)_d>\dim(A_D)_d$.

As in G-11, the only polynomial $\mathrm{GL}_4$
representation of weight size $nd$ that is trivial on
$\mathrm{SL}_4$ is the one-dimensional determinant power
of shape $(nd/4)^4$. If $4\nmid nd$, both invariant
pieces are zero. Thus the successful degree necessarily
satisfies $4\mid nd$, and the invariant dimension gap is
the desired rectangular multiplicity gap. The first-four-
variable highest-weight argument in G-05/R-14 transfers
it to the full orbit closures.

Finally $\beta=O(n^{33})$ and $u=O(n^2)$, so the explicit
formula \eqref{G13:first-bound} gives
$\log d_*=O(n^2\log n)$. The particular padded ranges
follow from R-14/R-16 and R-18, respectively; their
dimension gaps and $n<4m/3$ are the same hypotheses already
used to construct the stable padded point in G-11.
\end{proof}
